% ASSERT_CONVENTION: natural_units=natural, metric_signature=mostly_minus, fourier_convention=physics, coupling_convention=alpha_s, generator_normalization=T(fund)=1/2, covariant_derivative_sign=D_mu=partial_mu-igA
%
% Exercise Set 8: Physical Mesons as Dual Mesons
% Part of: The Dualised Standard Model -- Part III Exercises
%
% Conventions (Seiberg hep-th/9411149, unchanged from Lectures 1-10):
%   Quark flavour ordering: (u, d, s, c, b, t) = (1, 2, 3, 4, 5, 6) -- PDG standard
%   Meson naming: PDG 2024 naming conventions
%   Q_em(meson) = Q_em(q_i) - Q_em(q_j) for M_{ij} = q_i qbar_j
%   ChPT: U = exp(2i pi^a T^a / f_pi), f_pi = 92.1 MeV
%   Linear sigma model: phi = (v + sigma + i pi)/sqrt(2)
%   Epistemic: all non-SUSY duality claims carry conjectural qualifiers (P1)
%
\documentclass[11pt,a4paper]{article}
% ASSERT_CONVENTION: natural_units=natural, metric_signature=mostly_minus, fourier_convention=physics, coupling_convention=alpha_s, generator_normalization=T(fund)=1/2, covariant_derivative_sign=D_mu=partial_mu-igA
%
% CONVENTIONS (Seiberg hep-th/9411149)
% Metric: (+,-,-,-) (mostly minus)
% T(fund) = 1/2 per Weyl fermion
% b_0 = (11N_c - 2N_f)/3 for N_f Dirac flavours
% D_mu = partial_mu - ig A_mu^a T^a
% A(fund) = 1 (cubic anomaly coefficient per fundamental Weyl fermion)
% alpha_s = g_s^2 / (4 pi)
% R[Q] = (N_f - N_c)/N_f; R[W] = 2
% N_f counts Dirac flavour pairs (each = 2 Weyl fermions)
%
% This file is \input{} by all lecture files.

%% ---------- Document class ----------
% (loaded by the wrapper; preamble only provides packages and macros)

%% ---------- Standard packages ----------
\usepackage{amsmath,amssymb,amsthm,mathtools}
\usepackage{hyperref}
\usepackage[capitalise,noabbrev]{cleveref}
\usepackage{enumitem}
\usepackage{booktabs}
\usepackage{array}
\usepackage{xcolor}
\usepackage{tikz}
\usetikzlibrary{arrows.meta,decorations.markings,positioning,calc}
\usepackage{fancybox}
\usepackage{tcolorbox}
\tcbuselibrary{skins,breakable}
\usepackage{slashed}              % Feynman slash notation
\usepackage{cite}                 % compressed citation lists

%% ---------- Hyperref configuration ----------
\hypersetup{
  colorlinks=true,
  linkcolor=blue!70!black,
  citecolor=green!50!black,
  urlcolor=blue!80!black,
  pdfauthor={Part III Lecture Notes},
  pdftitle={The Dualised Standard Model}
}

%% ---------- Theorem environments ----------
\theoremstyle{plain}
\newtheorem{theorem}{Theorem}[section]
\newtheorem{lemma}[theorem]{Lemma}
\newtheorem{proposition}[theorem]{Proposition}
\newtheorem{corollary}[theorem]{Corollary}

\theoremstyle{definition}
\newtheorem{definition}[theorem]{Definition}
\newtheorem{example}[theorem]{Example}
\newtheorem{exercise}{Exercise}[section]

\theoremstyle{remark}
\newtheorem{remark}[theorem]{Remark}

%% ---------- Physics macros: gauge groups and representations ----------
\newcommand{\Nc}{N_{\!c}}
\newcommand{\Nf}{N_{\!f}}
\newcommand{\SU}[1]{\mathrm{SU}(#1)}
\newcommand{\UU}[1]{\mathrm{U}(#1)}

% Representations
\newcommand{\fund}{\square}
\newcommand{\antifund}{\overline{\square}}
\newcommand{\adj}{\mathrm{adj}}
\newcommand{\antisym}{\wedge^{\!2}}        % antisymmetric rank-2
\newcommand{\sym}{\mathrm{S}_2}           % symmetric rank-2

%% ---------- Physics macros: group theory constants ----------
% Dynkin index: Tr(T^a T^b) = T(R) delta^{ab}
\newcommand{\Tfund}{\tfrac{1}{2}}          % T(fund) = 1/2 per Weyl fermion
\newcommand{\Tadj}{\Nc}                    % T(adj) = N_c
\newcommand{\Cfund}{\frac{\Nc^2 - 1}{2\Nc}}  % C_2(fund)
\newcommand{\Cadj}{\Nc}                    % C_2(adj)

% Cubic anomaly coefficient: A(R) = Tr_R[T^a {T^b, T^c}] with A(fund) = 1
\newcommand{\Afund}{1}

%% ---------- Physics macros: beta function ----------
\newcommand{\bzero}{b_0}
\newcommand{\bone}{b_1}
\newcommand{\alphas}{\alpha_s}
\newcommand{\gs}{g_s}

%% ---------- Physics macros: SUSY / R-charge ----------
\newcommand{\Rcharge}[1]{R[#1]}
\newcommand{\Wsup}{W}                      % superpotential

%% ---------- Physics macros: covariant derivative and field strength ----------
\newcommand{\Dmu}{D_\mu}
\newcommand{\Fmn}{F_{\mu\nu}}
\newcommand{\Ftilde}{\widetilde{F}_{\mu\nu}}

%% ---------- Physics macros: common abbreviations ----------
\newcommand{\ie}{\textit{i.e.}}
\newcommand{\eg}{\textit{e.g.}}
\newcommand{\cf}{\textit{cf.}}
\newcommand{\IR}{\ensuremath{\mathrm{IR}}}
\newcommand{\UV}{\ensuremath{\mathrm{UV}}}
\newcommand{\AF}{\ensuremath{\mathrm{AF}}}
\newcommand{\BZ}{\ensuremath{\mathrm{BZ}}}
\newcommand{\NSVZ}{\ensuremath{\mathrm{NSVZ}}}
\newcommand{\SQCD}{\ensuremath{\mathrm{SQCD}}}
\newcommand{\QCD}{\ensuremath{\mathrm{QCD}}}
\newcommand{\SM}{\ensuremath{\mathrm{SM}}}
\newcommand{\Tr}{\mathrm{Tr}}

%% ---------- Convention box macro ----------
\newtcolorbox{conventionbox}[1][]{%
  enhanced,
  colback=blue!5!white,
  colframe=blue!60!black,
  fonttitle=\bfseries,
  title={Conventions},
  #1
}

%% ---------- Comparison table style ----------
\newtcolorbox{comparisontable}[1][]{%
  enhanced,
  colback=orange!5!white,
  colframe=orange!70!black,
  fonttitle=\bfseries,
  title={Comparison},
  #1
}

%% ---------- Warning / important box ----------
\newtcolorbox{warningbox}[1][]{%
  enhanced,
  colback=red!5!white,
  colframe=red!70!black,
  fonttitle=\bfseries,
  title={Important},
  #1
}

%% ---------- Preview / forward-reference box ----------
\newtcolorbox{previewbox}[1][]{%
  enhanced,
  colback=green!5!white,
  colframe=green!50!black,
  fonttitle=\bfseries,
  title={Preview},
  #1
}

%% ---------- Page layout ----------
\usepackage[margin=2.5cm]{geometry}
\setlength{\parskip}{0.4em}
\setlength{\parindent}{0em}

%% ---------- Section formatting ----------
\usepackage{titlesec}
\titleformat{\section}{\Large\bfseries}{\thesection}{1em}{}
\titleformat{\subsection}{\large\bfseries}{\thesubsection}{1em}{}

%% ---------- End of preamble ----------


%% ---------- Additional macros for this exercise ----------
\newcommand{\fpi}{f_\pi}
\newcommand{\Lchi}{\Lambda_\chi}
\newcommand{\qqbar}{\langle \bar{q}\,q \rangle}

\title{\textbf{Exercise Set 8: Physical Mesons as Dual Mesons}}
\author{The Dualised Standard Model --- Part III Exercises}
\date{Lent Term 2027}

\begin{document}
\maketitle

\begin{conventionbox}[title={Conventions for these exercises}]
\renewcommand{\arraystretch}{1.3}
\begin{tabular}{@{}l l@{}}
\toprule
\textbf{Quantity} & \textbf{Convention} \\
\midrule
Units & $\hbar = c = 1$ (natural) \\
Metric & $\eta_{\mu\nu} = \mathrm{diag}(+1,-1,-1,-1)$ (mostly minus) \\
Flavour ordering & $(u, d, s, c, b, t) = (1, 2, 3, 4, 5, 6)$ (PDG) \\
Meson naming & PDG 2024~\cite{PDG2024} \\
Pion decay constant & $\fpi = 92.1\;\text{MeV}$ (physical convention) \\
Epistemic & ``if the conjectured duality holds'' for all non-SUSY claims (P1) \\
\bottomrule
\end{tabular}

\medskip

\noindent References: Lecture~10 (all sections);
PDG 2024~\cite{PDG2024}; Cacciapaglia--Sannino~\cite{CacciapagliaSannino2024}.

\medskip

\noindent\textbf{Total marks: 65. \quad Suggested time: 90 minutes.}

\medskip

\noindent\textbf{Epistemic note.}
The quantum number assignments and the GMOR relation in these exercises
are standard QCD results, independent of any duality.  The
\emph{interpretation} of the meson matrix as an elementary field of the
magnetic theory is conjectural and is marked explicitly where it arises.
\end{conventionbox}

\bigskip


%% ========================================================================
%%  PROBLEM 1: Meson Quantum Numbers
%% ========================================================================
\section*{Problem 1: Meson Quantum Numbers from the Flavour Matrix
  \hfill $\star\star$ \normalsize(30 min, 25 marks)}
\addcontentsline{toc}{section}{Problem 1: Meson quantum numbers}
\label{prob:meson-qn}

\noindent\textbf{Background.}
The meson matrix $\mathcal{M}_{ij}$ ($i,j = 1,\ldots,6$) has entries
$\mathcal{M}_{ij} \sim q_i\,\bar{q}_j$, where $q_i$ is a quark of
flavour $i$ and $\bar{q}_j$ an antiquark of flavour $j$.  Each entry
corresponds to a meson whose electric charge is (Lecture~10, Eq.~(10)):
\begin{equation}\label{eq:meson-charge-ex}
  Q_{\text{em}}(\mathcal{M}_{ij})
  = Q_{\text{em}}(q_i) - Q_{\text{em}}(q_j)\,.
\end{equation}
The quark quantum numbers are:
\begin{equation}\label{eq:quark-qn-table}
\renewcommand{\arraystretch}{1.2}
\begin{array}{l|cccccc}
  & u & d & s & c & b & t \\
\hline
Q_{\text{em}} & +\tfrac{2}{3} & -\tfrac{1}{3} & -\tfrac{1}{3}
  & +\tfrac{2}{3} & -\tfrac{1}{3} & +\tfrac{2}{3} \\[3pt]
I_3 & +\tfrac{1}{2} & -\tfrac{1}{2} & 0 & 0 & 0 & 0 \\[3pt]
S & 0 & 0 & -1 & 0 & 0 & 0 \\
C & 0 & 0 & 0 & +1 & 0 & 0 \\
B_q & 0 & 0 & 0 & 0 & -1 & 0 \\
\end{array}
\end{equation}
For a meson with $B = 0$ (baryon number), the generalised
Gell-Mann--Nishijima formula gives (Lecture~10, Eq.~(9)):
\begin{equation}\label{eq:GMN-ex}
  Q_{\text{em}} = I_3 + \tfrac{1}{2}(S + C + B_q)\,.
\end{equation}

\medskip

\noindent\textbf{Questions.}
\begin{enumerate}[label=(\alph*)]

  \item \textbf{[5 marks]}
    Write out the $3 \times 3$ meson matrix for the light quarks
    $(u, d, s)$.  Label each entry with its quark content
    $q_i\,\bar{q}_j$.  Identify the off-diagonal entries with
    physical pseudoscalar mesons ($\pi^\pm$, $K^\pm$, $K^0$,
    $\bar{K}^0$).

  \item \textbf{[10 marks]}
    For each of the six off-diagonal mesons in the $3 \times 3$ block:
    \begin{enumerate}[label=(\roman*)]
      \item Determine the quark content.
      \item Assign the quantum numbers $I$, $I_3$, and $S$ of the
        meson (using $I_3^{\text{meson}} = I_{3}^{q_i} - I_{3}^{q_j}$
        and $S^{\text{meson}} = S^{q_i} - S^{q_j}$).
      \item Verify the electric charge using the Gell-Mann--Nishijima
        formula~\eqref{eq:GMN-ex}.
      \item Also verify the charge directly from
        $Q = Q(q_i) - Q(q_j)$.
    \end{enumerate}
    Present your results in a table.

  \item \textbf{[5 marks]}
    The three diagonal entries $\mathcal{M}_{11} = u\bar{u}$,
    $\mathcal{M}_{22} = d\bar{d}$, $\mathcal{M}_{33} = s\bar{s}$
    are flavour-neutral.  Explain why the physical states are not
    simply $u\bar{u}$, $d\bar{d}$, $s\bar{s}$ but rather
    $\pi^0$, $\eta$, and $\eta'$.  What symmetry principle
    determines the mixing?  Briefly explain the role of the
    $\UU{1}_A$ anomaly in making $\eta'$ heavy.

  \item \textbf{[5 marks]}
    Extend the meson matrix to $\Nf = 4$ by including the charm
    quark.  Identify the new entries $\mathcal{M}_{14}$,
    $\mathcal{M}_{24}$, $\mathcal{M}_{34}$, $\mathcal{M}_{41}$,
    $\mathcal{M}_{42}$, $\mathcal{M}_{43}$ with physical mesons
    (the $D$~meson family).  State their quark content, electric
    charge, and PDG mass (to the nearest MeV).

\end{enumerate}


\bigskip
\hrule
\bigskip


%% ========================================================================
%%  SOLUTION TO PROBLEM 1
%% ========================================================================
\begin{proof}[\textbf{Solution to Problem 1}]

\textbf{(a) The $3 \times 3$ light-quark meson matrix.}

\[
  \mathcal{M}_{ij} =
  \begin{pmatrix}
    u\bar{u} & u\bar{d} & u\bar{s} \\
    d\bar{u} & d\bar{d} & d\bar{s} \\
    s\bar{u} & s\bar{d} & s\bar{s}
  \end{pmatrix}
  =
  \begin{pmatrix}
    \pi^0/\eta/\eta' & \pi^+ & K^+ \\
    \pi^- & \pi^0/\eta/\eta' & K^0 \\
    K^- & \bar{K}^0 & \eta_s
  \end{pmatrix}
\]

The diagonal entries involve flavour mixing (discussed in part~(c)).
The off-diagonal entries each correspond to a single physical
pseudoscalar meson.

\bigskip

\textbf{(b) Quantum numbers of off-diagonal mesons.}

For each meson $\mathcal{M}_{ij} = q_i\,\bar{q}_j$, the quantum
numbers are obtained from:
\begin{align*}
  I_3^{\text{meson}} &= I_3(q_i) - I_3(q_j)\,, \\
  S^{\text{meson}} &= S(q_i) - S(q_j)
    = S(q_i) + S(\bar{q}_j)\,.
\end{align*}

\renewcommand{\arraystretch}{1.3}
\begin{center}
\begin{tabular}{l c c c c c c c}
\toprule
\textbf{Meson} & \textbf{Entry} & \textbf{Content}
  & $I$ & $I_3$ & $S$ & $Q$ (GMN) & $Q$ (direct) \\
\midrule
$\pi^+$ & $\mathcal{M}_{12}$ & $u\bar{d}$
  & $1$ & $+1$ & $0$ & $+1 + 0 = +1$ & $\tfrac{2}{3} + \tfrac{1}{3} = +1$ \\
$\pi^-$ & $\mathcal{M}_{21}$ & $d\bar{u}$
  & $1$ & $-1$ & $0$ & $-1 + 0 = -1$ & $-\tfrac{1}{3} - \tfrac{2}{3} = -1$ \\
$K^+$ & $\mathcal{M}_{13}$ & $u\bar{s}$
  & $\tfrac{1}{2}$ & $+\tfrac{1}{2}$ & $+1$
  & $+\tfrac{1}{2} + \tfrac{1}{2} = +1$
  & $\tfrac{2}{3} + \tfrac{1}{3} = +1$ \\
$K^0$ & $\mathcal{M}_{23}$ & $d\bar{s}$
  & $\tfrac{1}{2}$ & $-\tfrac{1}{2}$ & $+1$
  & $-\tfrac{1}{2} + \tfrac{1}{2} = 0$
  & $-\tfrac{1}{3} + \tfrac{1}{3} = 0$ \\
$K^-$ & $\mathcal{M}_{31}$ & $s\bar{u}$
  & $\tfrac{1}{2}$ & $-\tfrac{1}{2}$ & $-1$
  & $-\tfrac{1}{2} - \tfrac{1}{2} = -1$
  & $-\tfrac{1}{3} - \tfrac{2}{3} = -1$ \\
$\bar{K}^0$ & $\mathcal{M}_{32}$ & $s\bar{d}$
  & $\tfrac{1}{2}$ & $+\tfrac{1}{2}$ & $-1$
  & $+\tfrac{1}{2} - \tfrac{1}{2} = 0$
  & $-\tfrac{1}{3} + \tfrac{1}{3} = 0$ \\
\bottomrule
\end{tabular}
\end{center}

\textbf{Verification:}
\begin{itemize}
  \item For kaons: $S = +1$ for $\bar{s}$ (antiquark carries
    opposite strangeness to quark; $s$ has $S = -1$, so
    $\bar{s}$ has $S = +1$).  The Gell-Mann--Nishijima formula
    $Q = I_3 + S/2$ correctly gives $Q(K^+) = 1/2 + 1/2 = 1$.
    \checkmark
  \item All six charges agree between the two methods.  \checkmark
\end{itemize}

\bigskip

\textbf{(c) Diagonal mixing and the $\eta'$ problem.}

The diagonal entries $u\bar{u}$, $d\bar{d}$, $s\bar{s}$ all have
$Q = 0$, $I_3 = 0$, $S = 0$.  Therefore they can mix.  The physical
states are determined by the \textbf{flavour $\SU{3}$ representation
decomposition}:
\[
  \mathbf{3} \otimes \bar{\mathbf{3}} = \mathbf{8} \oplus \mathbf{1}\,.
\]

\begin{itemize}
  \item The \textbf{octet} ($\mathbf{8}$) contains $\pi^0$ and
    $\eta_8$:
    \[
      \pi^0 = \frac{u\bar{u} - d\bar{d}}{\sqrt{2}}\,,\qquad
      \eta_8 = \frac{u\bar{u} + d\bar{d} - 2s\bar{s}}{\sqrt{6}}\,.
    \]
  \item The \textbf{singlet} ($\mathbf{1}$) is:
    \[
      \eta_1 = \frac{u\bar{u} + d\bar{d} + s\bar{s}}{\sqrt{3}}\,.
    \]
\end{itemize}

If $\SU{3}$ flavour symmetry were exact, $\eta_8$ and $\eta_1$ would
be mass eigenstates.  In reality, $\SU{3}$ is broken by
$m_s \gg m_u, m_d$, so $\eta_8$ and $\eta_1$ mix to produce the
physical $\eta$ ($548\;\text{MeV}$) and $\eta'$ ($958\;\text{MeV}$).

The \textbf{$\UU{1}_A$ anomaly} plays a critical role: the axial
$\UU{1}_A$ symmetry is broken by instantons, giving the $\eta_1$ an
anomalous mass contribution.  Without this anomaly, the $\eta'$ would
be a ninth pseudo-Goldstone boson with mass comparable to the pion.
The resolution of this ``$\UU{1}$ problem'' was provided by
't~Hooft~\cite{tHooft1976}, Witten~\cite{Witten1979VenCre}, and
Veneziano~\cite{Veneziano1979}: the anomaly generates a mass
\[
  m_{\eta'}^2 \sim \frac{2\Nf}{f_\pi^2}\,\chi_{\text{top}}\,,
\]
where $\chi_{\text{top}}$ is the topological susceptibility of
pure Yang--Mills theory.  This pushes $m_{\eta'}$ to $\sim 958$~MeV,
much heavier than the octet pseudoscalars.

\bigskip

\textbf{(d) Extension to $\Nf = 4$: the $D$~meson family.}

The new entries involving charm quarks and light antiquarks:

\renewcommand{\arraystretch}{1.3}
\begin{center}
\begin{tabular}{l c c c l}
\toprule
\textbf{Meson} & \textbf{Entry} & \textbf{Quark content} & $Q$ &
  \textbf{Mass (MeV)} \\
\midrule
$D^0$ & $\mathcal{M}_{41}$ & $c\bar{u}$
  & $\tfrac{2}{3} - \tfrac{2}{3} = 0$ & $1864.8$ \\
$D^+$ & $\mathcal{M}_{42}$ & $c\bar{d}$
  & $\tfrac{2}{3} + \tfrac{1}{3} = +1$ & $1869.7$ \\
$D_s^+$ & $\mathcal{M}_{43}$ & $c\bar{s}$
  & $\tfrac{2}{3} + \tfrac{1}{3} = +1$ & $1968.3$ \\
$\bar{D}^0$ & $\mathcal{M}_{14}$ & $u\bar{c}$
  & $\tfrac{2}{3} - \tfrac{2}{3} = 0$ & $1864.8$ \\
$D^-$ & $\mathcal{M}_{24}$ & $d\bar{c}$
  & $-\tfrac{1}{3} - \tfrac{2}{3} = -1$ & $1869.7$ \\
$D_s^-$ & $\mathcal{M}_{34}$ & $s\bar{c}$
  & $-\tfrac{1}{3} - \tfrac{2}{3} = -1$ & $1968.3$ \\
\bottomrule
\end{tabular}
\end{center}

\textbf{Quantum number checks:}
\begin{itemize}
  \item $D^0 = c\bar{u}$: charm $C = +1$, no strangeness.
    $Q = I_3 + (C+S)/2 = 0 + 1/2 = 1/2$?  No---$D^0$ has
    $I_3 = -1/2$ (it is an isodoublet partner of $D^+$).

    Actually, the charm meson isospin assignments depend on
    convention.  For $c\bar{u}$:
    $Q = Q(c) - Q(u) = 2/3 - 2/3 = 0$.  \checkmark
  \item $D_s^+ = c\bar{s}$: $C = +1$, $S = +1$ ($\bar{s}$ contributes
    $+1$).  $Q = 2/3 + 1/3 = +1$.  \checkmark
\end{itemize}

Note that $\mathcal{M}_{ij}$ and $\mathcal{M}_{ji}$ are
charge conjugates: $D^0 = c\bar{u}$ and
$\bar{D}^0 = u\bar{c}$.

\end{proof}


\bigskip
\hrule
\bigskip


%% ========================================================================
%%  PROBLEM 2: Scalar vs Pseudoscalar Mesons
%% ========================================================================
\section*{Problem 2: The Scalar--Pseudoscalar Distinction
  \hfill $\star\star\star$ \normalsize(30 min, 20 marks)}
\addcontentsline{toc}{section}{Problem 2: Scalar vs pseudoscalar mesons}
\label{prob:scalar-pseudo}

\noindent\textbf{Background.}
In the SUSY dual theory, the meson $M_{ij}$ is a chiral superfield
whose lowest component is a complex scalar ($J^P = 0^+$).  Physical
pions, however, are pseudoscalars ($J^P = 0^-$).  The resolution
lies in the pattern of chiral symmetry breaking.

The lowest component $\phi_{ij}$ of the meson superfield can be
decomposed around the chiral condensate as (Lecture~10, Eq.~(15)):
\begin{equation}\label{eq:sigma-pi}
  \phi_{ij} = \frac{1}{\sqrt{2}}
  \bigl(v_{ij} + \sigma_{ij} + i\,\pi_{ij}\bigr)\,,
\end{equation}
where $v_{ij}$ is the vacuum expectation value, $\sigma_{ij}$ is
the scalar (amplitude) fluctuation ($J^P = 0^+$), and $\pi_{ij}$ is
the pseudoscalar (phase) fluctuation ($J^P = 0^-$).

\medskip

\noindent\textbf{Questions.}
\begin{enumerate}[label=(\alph*)]

  \item \textbf{[5 marks]}
    Identify the physical particles in each sector for $\Nf = 3$:
    \begin{enumerate}[label=(\roman*)]
      \item The scalar fluctuation $\sigma_{12}$ corresponds to which
        physical meson?  What is its mass?
      \item The pseudoscalar fluctuation $\pi_{12}$ corresponds to which
        physical meson?  What is its mass?
      \item Why is $m_\sigma \gg m_\pi$?  Give a symmetry-based
        explanation.
    \end{enumerate}

  \item \textbf{[10 marks]}
    The pseudoscalar mesons are pseudo-Goldstone bosons of chiral
    symmetry breaking
    $\SU{\Nf}_L \times \SU{\Nf}_R \to \SU{\Nf}_V$.
    \begin{enumerate}[label=(\roman*)]
      \item How many Goldstone bosons are there for $\Nf = 2$?
        Identify them with physical mesons.
      \item How many Goldstone bosons for $\Nf = 3$?  Identify them.
      \item The pions are massive ($m_\pi \neq 0$) despite being
        Goldstone bosons.  Explain why, and state what would happen
        to $m_\pi$ in the chiral limit $m_u, m_d \to 0$.
      \item The Gell-Mann--Oakes--Renner (GMOR)
        relation~\cite{GMOR1968} connects the pion mass to the
        quark masses and the chiral condensate:
        \begin{equation}\label{eq:gmor-ex}
          m_\pi^2\, \fpi^2 = -(m_u + m_d)\,\qqbar\,.
        \end{equation}
        Using the PDG/FLAG inputs~\cite{PDG2024,FLAG2024}:
        \begin{align*}
          \fpi &= 92.1\;\text{MeV}\,, \\
          \qqbar &\approx -(270\;\text{MeV})^3\,, \\
          m_u + m_d &\approx 8.5\;\text{MeV}
            \quad(\overline{\text{MS}},\; \mu = 2\;\text{GeV})\,,
        \end{align*}
        compute $m_\pi$ numerically and compare with the observed
        value $m_{\pi^\pm} = 139.6\;\text{MeV}$.
    \end{enumerate}

  \item \textbf{[5 marks]}
    Explain why the statement ``the pion IS the dual meson'' is
    misleading without qualification.  Your explanation should
    address:
    \begin{enumerate}[label=(\roman*)]
      \item What is the $J^P$ of the lowest component of the SUSY
        meson superfield?
      \item What is the $J^P$ of the pion?
      \item How does the pion arise from the decomposition~\eqref{eq:sigma-pi}?
      \item What is the \textbf{correct} qualified statement relating
        the dual meson to the pion?
    \end{enumerate}
    (\textit{All claims about the duality interpretation should carry
    the qualifier ``if the conjectured duality holds.''}  The QCD
    statements about chiral symmetry breaking are standard.)

\end{enumerate}


\bigskip
\hrule
\bigskip


%% ========================================================================
%%  SOLUTION TO PROBLEM 2
%% ========================================================================
\begin{proof}[\textbf{Solution to Problem 2}]

\textbf{(a) Physical identification of scalar and pseudoscalar modes.}

\textit{(i) Scalar mode $\sigma_{12}$.}

The entry $(1,2)$ corresponds to $u\bar{d}$.  The scalar
($J^P = 0^+$) fluctuation $\sigma_{12}$ is the
$a_0(980)^+$ meson, with mass $\sim 980\;\text{MeV}$.

\textit{(ii) Pseudoscalar mode $\pi_{12}$.}

The pseudoscalar ($J^P = 0^-$) fluctuation $\pi_{12}$ is
the $\pi^+$ pion, with mass $139.6\;\text{MeV}$.

\textit{(iii) Why $m_\sigma \gg m_\pi$.}

The pion is a \textbf{pseudo-Goldstone boson} of spontaneous chiral
symmetry breaking
$\SU{\Nf}_L \times \SU{\Nf}_R \to \SU{\Nf}_V$.  Goldstone's theorem
guarantees that it would be exactly massless if the symmetry were
exact; its mass arises only from the small explicit breaking by
quark masses ($m_u, m_d \sim 5\;\text{MeV} \ll \Lchi$).

The scalar meson $a_0$, by contrast, is \textbf{not} protected by any
Goldstone mechanism.  It is a massive excitation of the condensate whose
mass is set by the strong interaction scale $\sim \Lambda_\chi
\sim 1\;\text{GeV}$.  Hence $m_{a_0} \sim 1\;\text{GeV} \gg
m_\pi \sim 140\;\text{MeV}$.

\bigskip

\textbf{(b) Goldstone boson counting and the GMOR relation.}

\textit{(i) $\Nf = 2$.}

$\SU{2}_L \times \SU{2}_R \to \SU{2}_V$.  The number of broken
generators is $\Nf^2 - 1 = 3$.  The three Goldstone bosons are
$\pi^+$, $\pi^-$, $\pi^0$.

\textit{(ii) $\Nf = 3$.}

$\SU{3}_L \times \SU{3}_R \to \SU{3}_V$.  The number of broken
generators is $\Nf^2 - 1 = 8$.  The eight Goldstone bosons are
$\pi^+$, $\pi^-$, $\pi^0$, $K^+$, $K^-$, $K^0$, $\bar{K}^0$,
$\eta$.  (The $\eta'$ is not a Goldstone boson of
$\SU{3} \times \SU{3}$ breaking; it is the $\UU{1}_A$ partner
made heavy by instantons.)

\textit{(iii) Why pions are massive.}

The pions are \textbf{pseudo-Goldstone bosons}: the chiral symmetry
$\SU{\Nf}_L \times \SU{\Nf}_R$ is not exact but is \emph{explicitly}
broken by the quark mass terms $m_q\,\bar{q}\,q$ in the QCD
Lagrangian.  This explicit breaking gives the pions a mass
proportional to the quark masses.

In the \textbf{chiral limit} $m_u, m_d \to 0$, the explicit breaking
vanishes and $m_\pi \to 0$.  The pions become true Goldstone bosons.

\textit{(iv) GMOR numerical computation.}

\begin{align*}
  m_\pi^2
  &= \frac{(m_u + m_d)\,|\qqbar|}{\fpi^2} \\
  &= \frac{8.5\;\text{MeV} \times (270\;\text{MeV})^3}
    {(92.1\;\text{MeV})^2}\,.
\end{align*}

\textbf{Step 1:} $(270)^3 = 270 \times 270 \times 270 = 19{,}683{,}000\;\text{MeV}^3$.

\textbf{Step 2:} Numerator: $8.5 \times 19{,}683{,}000
= 167{,}306{,}000\;\text{MeV}^4 \approx 1.673 \times 10^8\;\text{MeV}^4$.

\textbf{Step 3:} Denominator: $(92.1)^2 = 8482\;\text{MeV}^2$.

\textbf{Step 4:}
\[
  m_\pi^2 = \frac{1.673 \times 10^8}{8482}\;\text{MeV}^2
  \approx 19{,}720\;\text{MeV}^2\,.
\]

\textbf{Step 5:}
\[
  m_\pi = \sqrt{19{,}720}\;\text{MeV}
  \approx 140\;\text{MeV}\,.
\]
\[
  \boxed{m_\pi \approx 140\;\text{MeV}\,.}
\]

This matches the measured value $m_{\pi^\pm} = 139.6\;\text{MeV}$ to
better than $1\%$, well within the $\sim 10\%$ uncertainties on the
input parameters $m_u + m_d$ and $\qqbar$.
\quad$\checkmark$

\bigskip

\textbf{(c) Why ``the pion IS the dual meson'' is misleading.}

\textit{(i)} The lowest component $\phi_{ij}$ of the SUSY chiral
superfield $M_{ij}$ has $J^P = 0^+$ (scalar).

\textit{(ii)} The pion has $J^P = 0^-$ (pseudoscalar).

\textit{(iii)} The pion arises as the \textbf{imaginary part}
(phase fluctuation) of the decomposition $\phi_{ij} = (v_{ij} +
\sigma_{ij} + i\,\pi_{ij})/\sqrt{2}$ around the chiral condensate.
It is not the scalar field $\phi_{ij}$ itself, but a fluctuation of
$\phi_{ij}$ around its VEV.

\textit{(iv)} The \textbf{correct qualified statement} is:

\begin{quote}
\emph{If the conjectured non-SUSY duality holds, the pseudoscalar
meson $\pi^+$ corresponds to the phase fluctuation of the chiral
condensate $\langle \phi_{12} \rangle$ of the dual meson field.  The
scalar lowest component $\phi_{12}$ of the chiral superfield $M_{12}$
maps instead to the scalar meson $a_0(980)^+$.  The SUSY meson
superfield contains both sectors; the physical distinction arises
from chiral symmetry breaking.}
\end{quote}

The key point is that a single complex scalar field $\phi_{ij}$
yields \textbf{two} physical sectors after symmetry breaking: the
scalars $\sigma$ and the pseudoscalars $\pi$.

\end{proof}


\bigskip
\hrule
\bigskip


%% ========================================================================
%%  PROBLEM 3: Mass Hierarchy from VEV Structure
%% ========================================================================
\section*{Problem 3: Mass Hierarchy from the VEV Structure
  \hfill $\star\star$ \normalsize(30 min, 20 marks)}
\addcontentsline{toc}{section}{Problem 3: Mass hierarchy from VEV structure}
\label{prob:mass-hierarchy}

\noindent\textbf{Background.}
In the dualised Standard Model (Lectures~7--8), \emph{if the
conjectured duality holds}, the effective Yukawa couplings are
(Lecture~7, Eq.~(14)):
\begin{equation}\label{eq:eff-yukawa}
  Y_{ij} = y\,\frac{\langle H_{ij} \rangle}{v}\,,
\end{equation}
where $y$ is a single universal coupling and $\langle H_{ij} \rangle$
is the VEV of the $(i,j)$-th Higgs doublet.  The fermion masses are
therefore:
\begin{equation}\label{eq:fermion-mass}
  m_{ij}^{q} = y\,\langle H_{ij}^{q} \rangle\,,
  \qquad q \in \{u, d\}\,.
\end{equation}
The physical quark masses (diagonal in the mass basis) arise from
diagonalising $m_{ij}$, with the CKM matrix encoding the mismatch
between the up-type and down-type diagonalisations.

The VEV hierarchy suggested by Cacciapaglia and
Sannino~\cite{CacciapagliaSannino2024} is:
\begin{equation}\label{eq:VEV-hierarchy}
  \langle H_{33} \rangle \gg \langle H_{22} \rangle
  \gg \langle H_{11} \rangle\,,
\end{equation}
arising from higher-dimensional (Planck-suppressed) operators in the
electric theory with $\mathcal{O}(1)$ dimensionless coefficients.

\medskip

\noindent\textbf{Questions.}
\begin{enumerate}[label=(\alph*)]

  \item \textbf{[5 marks]}
    If $\langle H_{33}^u \rangle$ is the dominant VEV (comparable to
    the electroweak scale $v$), and the universal Yukawa coupling $y$
    is $\mathcal{O}(1)$:
    \begin{enumerate}[label=(\roman*)]
      \item Show that the top quark mass is
        $m_t = y\,\langle H_{33}^u \rangle \sim v \sim 174\;\text{GeV}$.
      \item Explain why this is a parametric consistency check
        (not a prediction): what are the free parameters?
    \end{enumerate}

  \item \textbf{[10 marks]}
    The observed quark masses at $\mu = 2\;\text{GeV}$ in the
    $\overline{\text{MS}}$ scheme are approximately~\cite{PDG2024}:
    \begin{equation}\label{eq:quark-masses}
    \renewcommand{\arraystretch}{1.2}
    \begin{array}{lrl}
      m_u \approx 2.2\;\text{MeV}\,, & \qquad
      m_c \approx 1270\;\text{MeV}\,, & \qquad
      m_t \approx 173{,}000\;\text{MeV}\,, \\
      m_d \approx 4.7\;\text{MeV}\,, &
      m_s \approx 95\;\text{MeV}\,, &
      m_b \approx 4180\;\text{MeV}\,.
    \end{array}
    \end{equation}
    (Note: $m_t$ is the pole mass; the others are running masses.)

    \begin{enumerate}[label=(\roman*)]
      \item Compute the quark mass ratios $m_t/m_c$, $m_c/m_u$,
        $m_b/m_s$, $m_s/m_d$ from the values above.
      \item If the VEV hierarchy is geometric:
        $\langle H_{33}\rangle/\langle H_{22}\rangle \sim 1/\epsilon$,
        $\langle H_{22}\rangle/\langle H_{11}\rangle \sim 1/\epsilon$,
        determine the value of $\epsilon$ needed to reproduce each of
        the ratios in~(i) (treating up-type and down-type sectors
        independently).
      \item Can a \textbf{single} value of $\epsilon$ simultaneously
        reproduce all four ratios?  What does this tell you about
        the $\mathcal{O}(1)$ coefficients?
    \end{enumerate}

  \item \textbf{[5 marks]}
    In the light-quark sector, the pseudoscalar meson masses
    ($m_\pi \approx 140\;\text{MeV}$,
    $m_K \approx 494\;\text{MeV}$,
    $m_\eta \approx 548\;\text{MeV}$) are determined by QCD
    dynamics (quark masses + chiral symmetry breaking), not
    directly by the Higgs VEVs.
    \begin{enumerate}[label=(\roman*)]
      \item Explain why the dualised SM framework is
        \textbf{consistent with}, but does not \textbf{improve upon},
        standard QCD predictions for light meson masses.
      \item What \textbf{does} the dualised SM add to the picture
        beyond standard QCD?  List at least two structural insights
        (assuming the conjectured duality holds).
    \end{enumerate}

\end{enumerate}


\bigskip
\hrule
\bigskip


%% ========================================================================
%%  SOLUTION TO PROBLEM 3
%% ========================================================================
\begin{proof}[\textbf{Solution to Problem 3}]

\textbf{(a) Top quark mass.}

\textit{(i)} With the dominant VEV $\langle H_{33}^u \rangle \sim v$
(the electroweak scale) and $y \sim \mathcal{O}(1)$:
\[
  m_t = y\,\langle H_{33}^u \rangle \sim y \times v
  \sim 1 \times 174\;\text{GeV} = 174\;\text{GeV}\,.
\]
(Here $v = (\sqrt{2}\,G_F)^{-1/2} \approx 246\;\text{GeV}$ and
$m_t = y_t\,v/\sqrt{2} \approx y_t \times 174\;\text{GeV}$.  Since
$y_t = y\,\langle H_{33}^u\rangle/v$ and $\langle H_{33}^u\rangle
\sim v$, we get $y_t \sim y \sim \mathcal{O}(1)$.)

This matches the observed $m_t \approx 173\;\text{GeV}$.
\quad$\checkmark$

\textit{(ii)} This is a \textbf{parametric consistency check}, not a
prediction, because:
\begin{itemize}
  \item The coupling $y$ is a free parameter (it is $\mathcal{O}(1)$
    by assumption, not predicted to be a specific number).
  \item The VEV $\langle H_{33}^u \rangle$ is set by the minimisation
    of the scalar potential, which depends on unknown
    $\mathcal{O}(1)$ coefficients from Planck-suppressed operators.
  \item The statement ``$m_t \sim v$'' is equivalent to ``$y_t \sim
    \mathcal{O}(1)$'', which is an input assumption, not an output.
\end{itemize}

The non-trivial content is that the framework \emph{naturally
accommodates} $m_t \sim v$ with $\mathcal{O}(1)$ couplings, without
requiring fine-tuning.

\bigskip

\textbf{(b) Quark mass ratios and the geometric hierarchy.}

\textit{(i) Mass ratios.}

\renewcommand{\arraystretch}{1.3}
\begin{center}
\begin{tabular}{l r}
\toprule
\textbf{Ratio} & \textbf{Value} \\
\midrule
$m_t / m_c$ & $173{,}000 / 1270 \approx 136$ \\
$m_c / m_u$ & $1270 / 2.2 \approx 577$ \\
$m_b / m_s$ & $4180 / 95 \approx 44$ \\
$m_s / m_d$ & $95 / 4.7 \approx 20$ \\
\bottomrule
\end{tabular}
\end{center}

\textit{(ii) Geometric hierarchy parameter.}

In the geometric hierarchy picture, the mass of the $n$-th generation
scales as $m_n \propto \epsilon^{3-n}$ (with the 3rd generation
having the largest mass).  The ratio between consecutive generations is
$m_{n+1}/m_n \sim 1/\epsilon$.  Therefore:

\textbf{Up-type sector:}
\[
  \epsilon_u^{(t/c)} = \frac{m_c}{m_t}
  = \frac{1}{136} \approx 0.0074\,,
  \qquad
  \epsilon_u^{(c/u)} = \frac{m_u}{m_c}
  = \frac{1}{577} \approx 0.0017\,.
\]

\textbf{Down-type sector:}
\[
  \epsilon_d^{(b/s)} = \frac{m_s}{m_b}
  = \frac{1}{44} \approx 0.023\,,
  \qquad
  \epsilon_d^{(s/d)} = \frac{m_d}{m_s}
  = \frac{1}{20} = 0.050\,.
\]

\textit{(iii) Can a single $\epsilon$ work?}

\textbf{No.}  The four values of $\epsilon$ range from $0.0017$ to
$0.050$---a factor of $\sim 30$.  A \textbf{single} geometric ratio
$\epsilon$ cannot simultaneously reproduce all four mass ratios.

This is \textbf{expected} in the Cacciapaglia--Sannino
framework~\cite{CacciapagliaSannino2024}: the VEV hierarchy is
generated by higher-dimensional operators with $\mathcal{O}(1)$
dimensionless coefficients $c^{abcd}$.  These coefficients are
different for each entry of the Higgs matrix, so the effective
$\epsilon$ varies between sectors:
\[
  m_i \sim y\,c_i\,\epsilon^{3-n_i}\,v\,,
\]
where $c_i$ are $\mathcal{O}(1)$ constants (not predicted by the
framework).

The \textbf{parametric consistency check} is that
$\epsilon \in [0.002, 0.05]$---spanning a factor of $25$---can be
absorbed into $\mathcal{O}(1)$ coefficients.  Specifically, if
$\epsilon \sim 0.01$ is a ``typical'' suppression per generation, then
the actual ratios $(0.002 \text{ to } 0.05)$ differ from $0.01$ by
factors of $0.2$ to $5$, which are all $\mathcal{O}(1)$.

\[
  \boxed{\text{Parametric consistency: all ratios explained by a
  single } \epsilon \sim 0.01 \text{ with } \mathcal{O}(1)
  \text{ coefficients.}}
\]

\bigskip

\textbf{(c) Light meson masses and the dualised SM.}

\textit{(i)} The dualised SM framework is \textbf{consistent with}
standard QCD for light meson masses because:
\begin{itemize}
  \item The pseudoscalar meson masses ($m_\pi$, $m_K$, $m_\eta$)
    are determined by chiral symmetry breaking, which depends on
    the quark masses $m_u$, $m_d$, $m_s$ and the chiral condensate
    $\qqbar$.
  \item The quark masses enter through the VEV hierarchy (via
    $m_q = y\,\langle H \rangle$), but the meson masses depend on
    $m_q$ \emph{and} on the non-perturbative QCD dynamics that
    form the condensate.
  \item The dualised SM does not predict $\qqbar$ or the
    individual quark masses; it only provides the hierarchical
    \emph{structure} $m_t \gg m_c \gg m_u$ with $\mathcal{O}(1)$
    free coefficients.
\end{itemize}

Therefore, the dualised SM reproduces the standard QCD predictions for
light meson masses, but does not improve upon them: any set of quark
masses can be accommodated by adjusting the $\mathcal{O}(1)$
coefficients.

\textit{(ii)} What the dualised SM adds (assuming the conjectured
duality holds):

\begin{enumerate}
  \item \textbf{A UV origin for the meson matrix structure:}  The
    $6 \times 6$ meson matrix arises as an elementary field in the
    magnetic theory, rather than being an ad hoc effective field.
    The $\Nf \times \Nf$ matrix structure is dictated by the gauge
    group of the electric theory.

  \item \textbf{The universal Yukawa coupling:} All Yukawa couplings
    arise from a single coupling $y$ in the magnetic superpotential,
    with the observed hierarchy coming from the VEV structure rather
    than from a hierarchy of couplings.  This replaces 13 free Yukawa
    parameters (in the SM) with 1 coupling + Planck-scale operators
    (though the operators introduce their own free coefficients).

  \item \textbf{Connection between the Higgs sector and strong
    dynamics:}  The Higgs doublets are identified with entries of the
    meson matrix, providing a structural link between
    electroweak symmetry breaking and the strong interaction sector.
\end{enumerate}

\end{proof}


\bigskip

\begin{center}
\rule{0.6\textwidth}{0.4pt}

\textit{End of Exercise Set 8.  These exercises illustrate three
aspects of the physical meson identification: the quantum numbers of
the flavour meson matrix are completely determined by quark content
and the Gell-Mann--Nishijima formula (Problem~1); the distinction
between scalar and pseudoscalar mesons arises from chiral symmetry
breaking, with the GMOR relation confirming
$m_\pi \approx 140\;\text{MeV}$ (Problem~2); and the fermion mass
hierarchy can be parametrically accommodated by a VEV hierarchy with
$\mathcal{O}(1)$ coefficients, consistent with but not predictive
beyond standard QCD (Problem~3).}
\end{center}


%% ========================================================================
%%  BIBLIOGRAPHY
%% ========================================================================
\begin{thebibliography}{99}

\bibitem{PDG2024}
  R.~L.~Workman \textit{et al.} (Particle Data Group),
  ``Review of Particle Physics,''
  \textit{Phys.\ Rev.\ D} \textbf{110} (2024) 030001.

\bibitem{FLAG2024}
  Y.~Aoki \textit{et al.} (FLAG),
  ``FLAG Review 2024,''
  \textit{Eur.\ Phys.\ J.\ C} \textbf{84} (2024) 148
  [\texttt{arXiv:2111.09849}].

\bibitem{CacciapagliaSannino2024}
  G.~Cacciapaglia and F.~Sannino,
  ``The Dualised Standard Model,''
  \texttt{arXiv:2407.17281} [hep-ph] (2024).

\bibitem{GMOR1968}
  M.~Gell-Mann, R.~J.~Oakes, and B.~Renner,
  ``Behavior of current divergences under $\SU{3} \times \SU{3}$,''
  \textit{Phys.\ Rev.} \textbf{175} (1968) 2195.

\bibitem{tHooft1976}
  G.~'t~Hooft,
  ``Computation of the quantum effects due to a four-dimensional
  pseudoparticle,''
  \textit{Phys.\ Rev.\ D} \textbf{14} (1976) 3432.

\bibitem{Witten1979VenCre}
  E.~Witten,
  ``Current algebra theorems for the $\UU{1}$ `Goldstone boson',''
  \textit{Nucl.\ Phys.} \textbf{B156} (1979) 269.

\bibitem{Veneziano1979}
  G.~Veneziano,
  ``$\UU{1}$ without instantons,''
  \textit{Nucl.\ Phys.} \textbf{B159} (1979) 213.

\bibitem{Seiberg1995}
  N.~Seiberg,
  ``Electric--magnetic duality in supersymmetric non-Abelian gauge
  theories,''
  \textit{Nucl.\ Phys.} \textbf{B435} (1995) 129
  [\texttt{hep-th/9411149}].

\end{thebibliography}

\end{document}
